\section{Results}
\label{sec:results}
Scores are mSCF1 at $K_{\text{bin}}$, seed~1, bin level, paired by section,
unless stated otherwise. Labels are quoted exactly.

\subsection{Baseline reproductions (RQ4)}
\label{sec:res-repro}
Each baseline was run under its own protocol and evaluator.
\begin{itemize}
  \item \textbf{msiPL} reproduced on all three collections: GBM 0.414 against
        0.403 published, CAC 0.402 against 0.431, and RCC 0.157 against 0.159.
  \item \textbf{S3PL} reproduced on CAC (0.592 against 0.601) and on RCC
        (0.41--0.44 over three seeds, against 0.47).
  \item \textbf{S3PL on GBM was not reproduced:} 0.32 against 0.496, using the
        current public code and released settings. The paper's supplementary
        settings were not available (\cref{sec:res-s3pl}).
\end{itemize}

\subsection{Correlation outperforms IG (RQ1, RQ2)}
\label{sec:res-main}
\Cref{tab:main} scores every ranking on the same segmentation and budget.

\textbf{Against msiPL.} Both VAE--IG arms beat reproduced msiPL in every GBM
and CAC section: 0.518 and 0.511 against 0.414 on GBM, and 0.528 and 0.560
against 0.402 on CAC. This is statistically supported in all four comparisons
(HL $+0.09$ to $+0.16$, Holm $p=0.031$). It compares whole pipelines, and one
predeclared clause (budget stability) could not be assessed.

\textbf{IG against correlation.} With the GMM held fixed, correlation with the
posterior maps beat IG in all 47 valid comparisons at seed~1
(\cref{fig:igcorr}). The predeclared rule's exact fall-through wording
applies: ``neither predeclared verdict applies, because IG scored below
posterior-$|$PCC$|$ in every section''. Over seeds 1--3, correlation won 136
of 140 comparisons (\cref{tab:seeds}).
\begin{itemize}
  \item The strict per-seed rule failed only at CAC seed~2, where IG was below
        in 6/8 sections instead of the required 7/8, so the label is ``does not
        hold overall''.
  \item No collection mean ever reversed direction.
  \item With seeds averaged within section, the difference is statistically
        supported in all six dataset--arm comparisons (\cref{tab:seeds}). Every
        section favours correlation ($r=-1$), and every raw $p$ is the
        smallest attainable for its $n$.
  \item The mixed models agree ($-0.24$ [$-0.27$, $-0.22$] on GBM, $-0.14$
        [$-0.17$, $-0.12$] on RCC) and attribute almost no variance to the seed.
  \item GBM is negative in every PDX line.
  \item Sampling 48 or 192 pixels for IG instead of 12 raised scores by at most
        0.024.
\end{itemize}
\emph{Answer:} the pipeline beats msiPL, but not because of IG. Given the same
segmentation, direct correlation is better on all three datasets. The useful
signal lies in the label-free segmentation.

\begin{table*}[t]
\caption{Collection mean mSCF1 at $K_{\text{bin}}$ (seed~1, bin level). All
rankings except the baselines use the same GMM. Bold: best label-free ranking.
The oracle uses the masks and is a ceiling.}
\label{tab:main}
\centering\small
\begin{tabular}{@{}llccccccccc@{}}
\toprule
Coll. & VAE arm & IG & Post.\ corr. & Moran's $I$ & Mean int. & Variance & msiPL & Random & Oracle \\
\midrule
GBM & centre & 0.518 & \textbf{0.743} & 0.684 & 0.483 & 0.497 & 0.414 & 0.009 & 0.809 \\
GBM & uniform & 0.511 & \textbf{0.744} & 0.684 & 0.483 & 0.497 & 0.414 & 0.009 & 0.809 \\
CAC & centre & 0.528 & 0.606 & 0.605 & \textbf{0.612} & 0.593 & 0.402 & 0.108 & 0.654 \\
CAC & uniform & 0.560 & 0.607 & 0.605 & \textbf{0.612} & 0.593 & 0.402 & 0.108 & 0.654 \\
RCC & centre & 0.412 & 0.568 & \textbf{0.631} & 0.498 & 0.508 & --- & 0.171 & 0.756 \\
RCC & uniform$^{a}$ & 0.410 & 0.596 & \textbf{0.635} & 0.499 & 0.504 & --- & 0.173 & 0.755 \\
\bottomrule
\multicolumn{10}{@{}l}{\footnotesize $^{a}$7 of 8 runs (UH0505\_12 failed the predeclared IG completeness check). msiPL on RCC is not budget-matched.}
\end{tabular}
\end{table*}

\begin{table}[t]
\caption{IG $-$ posterior correlation across seeds 1--3 (bin level): mean
$\pm$ SD of the seed-level collection means, wins for correlation, and paired
statistics on seed-averaged sections (HL with exact 95\% CI; Holm within F1).}
\label{tab:seeds}
\centering\scriptsize
\setlength{\tabcolsep}{3pt}
\begin{tabular}{@{}lcccc@{}}
\toprule
 & Mean $\pm$ SD & Corr.\ wins & HL [95\% CI] & Holm $p$ \\
\midrule
GBM centre & $-0.244 \pm 0.018$ & 24/24 & $-0.248$ [$-0.293$, $-0.186$] & 0.047 \\
GBM uniform & $-0.245 \pm 0.011$ & 23/23 & $-0.243$ [$-0.314$, $-0.159$] & 0.047 \\
CAC centre & $-0.077 \pm 0.003$ & 23/24 & $-0.078$ [$-0.116$, $-0.040$] & 0.047 \\
CAC uniform & $-0.055 \pm 0.007$ & 22/24 & $-0.046$ [$-0.080$, $-0.027$] & 0.047 \\
RCC centre & $-0.144 \pm 0.011$ & 22/23 & $-0.147$ [$-0.207$, $-0.082$] & 0.047 \\
RCC uniform & $-0.141 \pm 0.041$ & 22/22 & $-0.131$ [$-0.204$, $-0.056$] & 0.047 \\
\bottomrule
\end{tabular}
\end{table}

\begin{figure}[t]
\centering
\includegraphics[width=\columnwidth]{05_ig_minus_correlation_by_section.png}
\caption{IG minus posterior correlation per section and VAE arm (seed~1). All
47 valid comparisons are negative.}
\label{fig:igcorr}
\end{figure}

\subsection{Simple rankings and the metric}
\label{sec:res-simple}
The best predeclared label-free ranking differs by dataset:
\begin{itemize}
  \item GBM: posterior correlation (0.743);
  \item CAC: mean intensity (0.612), with correlation, Moran's~$I$ and variance
        all within 0.02;
  \item RCC: Moran's~$I$ (0.63).
\end{itemize}
CAC is near saturation: every method, including the oracle, converges near
0.41 at twice the budget. On GBM and RCC, sections whose GMM matches the masks
better have correlation scores closer to the oracle (Spearman $\rho=-0.90$ and
$-0.84$). mSCF1 rewards ion images that follow the masks, so mask-like
segmentations score well when correlated against.

\subsection{Spatial context (RQ3)}
\label{sec:res-spatial}
\begin{itemize}
  \item \textbf{Uniform $3\times3$ context} ``helps'' IG on CAC ($+0.023$ over
        three seeds, 8/8 sections), but it is not statistically supported
        (Holm $p=0.13$) and is 0.000 under correlation ranking. On GBM the
        result is ``incomplete (7/8)'', with a mean of $+0.000$.
  \item \textbf{It is not spatially specific.} On CAC, real and shuffled
        nonlocal neighbours tied (0.542 against 0.543). On GBM, zeroed context
        scored highest.
  \item \textbf{Larger windows} ($7\times7$, $9\times9$) were $0.035$ below
        $3\times3$ on GBM (``mixed'') and ``equivalent'' on CAC.
  \item \textbf{Learned attention} was ``equivalent'' to uniform averaging
        ($+0.011$ GBM, $-0.006$ CAC). Its weights stayed near uniform with real
        neighbours (normalised entropy 0.98--0.999) but became selective with
        shuffled ones. The mechanism works; real neighbours are simply
        near-copies of the centre (median cosine about 0.98). Centre masking
        did not make attention selective and lowered GBM scores by 0.039.
  \item \textbf{Smoothed segmentation (Test~S)} was never labelled ``helps'':
        four ``equivalent'', one ``does not meet the `helps' rule (labelled
        mixed)'', and one ``incomplete''.
\end{itemize}
\emph{Answer:} spatial input to the encoder or the segmentation gave no
consistent, spatially specific gain. Spatial information helped as a ranking
statistic, through correlation with coherent segments and Moran's~$I$.

\subsection{Is the nonlinear encoder needed? (RQ2)}
\label{sec:res-encoder}
Replacing the VAE with PCA-5 while keeping the GMM and correlation gave VAE
$-$ PCA differences of $+0.023$ on GBM (5/8, ``mixed''), $-0.013$ on CAC
(``mixed'') and ``equivalent'' on RCC. These were unchanged with the VAE
averaged over three seeds. The VAE segmented slightly better (balanced
accuracy 0.868 against 0.839 on GBM) without reliably better peaks.

PCA-5 also had \emph{lower} reconstruction error than the VAE on all eight GBM
sections, and higher error on 7/8 CAC sections, so reconstruction did not
track peak quality~\citep{mazarura2025autoencoder}. BIC never validated the
number of components: it chose the largest candidate ($K=6$) every time.

\subsection{S3PL: neighbours, reproduction and seeds (RQ3, RQ4)}
\label{sec:res-s3pl}
\textbf{Neighbour dependence.} On CAC, retraining S3PL with every patch filled
with the centre spectrum collapsed it from 0.591 to 0.18 and 0.15 (two seeds,
8/8 sections). Shuffled nonlocal neighbours recovered most of the score (0.55
against 0.59), so correct placement added only about $+0.04$, inconsistently
(``mixed'', Holm $p=0.63$). Training for longer made S3PL worse: epoch 10 was
best in 15/16 sections.

\textbf{Reproduction gap.} Over ten released seeds (\cref{fig:s3plseeds}),
S3PL scored $0.319\pm0.042$ on GBM and $0.547\pm0.038$ on CAC.
\begin{itemize}
  \item \textbf{GBM:} the published 0.496 is about 4 SD above our mean. It is
        above every one of our seeds in 6/8 sections, and best-of-ten per
        section reaches only 0.434. Labels: ``outside our seed distribution''
        and ``not explained by seed selection alone''.
  \item \textbf{CAC:} the published 0.601 is just above our maximum (0.591),
        within its two-decimal rounding of best-of-ten (0.5998).
  \item \textbf{Normalisation:} switching to the TIC normalisation described in
        the paper did not raise GBM (0.303 against 0.324; HL $-0.012$
        [$-0.154$, $0.099$]), so it ``does not explain'' the gap.
\end{itemize}

\textbf{Why seeds matter.} About 52\% of the score variance is between
sections, 0--1\% is a seed main effect, and 47\% is a section $\times$ seed
interaction.
\begin{itemize}
  \item Seeds switch particular sections between distinct outcomes:
        GBM108\_positive moves between 0.02 and 0.46, and CAC 520TopL between
        0.17 and 0.71.
  \item The peak lists of different seeds can overlap by as little as 1\%.
  \item Varying only the weight initialisation, or only the data order, showed
        that initialisation dominated (SD 0.222 against 0.001 on
        GBM108\_positive), although order also mattered on GBM108\_negative.
  \item Same-seed reruns were bit-identical.
  \item Under TIC the per-section seed SD fell from 0.061 to 0.011.
\end{itemize}
\emph{Answer:} msiPL reproduces. S3PL reproduces except on GBM, where neither
seeds nor normalisation explain the gap. Its starting weights select between
distinct solutions, so single-seed comparisons against it are unreliable.

\begin{figure}[t]
\centering
\includegraphics[width=\columnwidth]{s3pl_released_seed_distributions.pdf}
\caption{S3PL mSCF1 per section over ten released seeds (dots), the seed mean
(bar) and the published value (diamond), using S3PL's own evaluator.}
\label{fig:s3plseeds}
\end{figure}

\subsection{RCC and runtime}
\label{sec:res-rcc}
\label{sec:res-runtime}
\textbf{RCC.} On the held-out RCC data, correlation beat IG in every valid run
(0/8 and 0/7 positive), and Moran's~$I$ was the best label-free ranking. When
only cancer and normal pixels are scored, the advantage shrinks: correlation
beats IG in 5/8 runs (centre-only) and 6/7 (uniform). The three-class result
therefore partly reflects separating tissue from the unannotated region.

\textbf{Runtime.} Given a segmentation, correlation took 1.1\,s per GBM section
against 770\,s for IG (CAC 0.03 against 6.0\,s; RCC 0.11 against 372\,s), and
every timed ranking reproduced the saved ranking exactly.
\begin{itemize}
  \item Training: cached VAE runs took a median of 32--35\,min (GBM), 1.3\,min
        (CAC) and 2.2\,min (RCC), against 13.7, 0.9 and 1.6\,min for 10-epoch
        S3PL.
  \item Complete workflow, as the proposal planned to measure: on 160TopL the
        VAE took 122--135\,s against 151--174\,s for S3PL, with 1.2 against
        1.3\,GiB of RAM and 237 against 315\,MiB of GPU memory.
\end{itemize}
These are native-protocol measurements, not an equal-work comparison.
